\section{Theorem targets and backward reductions}
\label{sec:theorem-targets}

\subsection{The largest symmetric permanent level space through the all-ones matrix}
\label{sec:permanent-target}

\paragraph{Open target.}
Let \(J_4\) be the all-ones matrix. Prove that
\[
 \max\{\dim_{\mathbb C}L : L\subseteq\operatorname{Sym}_4(\mathbb C)
   \text{ a linear subspace},\
   \operatorname{per}(J_4+B)=24\text{ for every }B\in L\}=2 .
\]
The first-row-and-column construction in Section~\ref{sec:gct} attains
two. Every larger space contains a three-plane. Thus the upper bound
asks for the exclusion of every independent symmetric triple
\(B_1,B_2,B_3\). In \(\operatorname{per}(J_4+\sum_i t_iB_i)\)
the coefficient of \(t_i\) is \(6\sum_{r,c}(B_i)_{rc}\).
For every triple with zero entry sum in each direction it therefore
suffices to exhibit one nonzero coefficient of degree two, three, or
four. The constant is \(24\) and the total degree is at most four.
There are \(6+10+15=31\) such coefficients. Lean checks the linear
formula, the equivalence over \(\mathbb C\), and the count of 31 indices.

\paragraph{A sufficient certificate at five.}
Let \(P=\mathbb F_5[y_1,\ldots,y_{18}]\), let
\(I_\star=(f_1,\ldots,f_{31})\) be the coefficient ideal in the
normalized star chart of Section~\ref{sec:gct-path}, and let \(D\in P\)
be the determinant of its corrected diagonal readout. The chart covers
the geometric special fiber. By the proper specialization argument in
that section, \(1\in I_\star\) would exclude complex three-planes.
A point in characteristic five would leave its lift to characteristic
zero to be determined. Over \(\mathbb C\), write the integral star
equations as \(f_j^{\mathbb Z}(e,z)=0\), where \(e\) and \(z\) are
the free edge and diagonal entries. The target is equivalent to
emptiness of the projection
\[
 E_{\mathrm{lift}}=\{e\in\mathbb C^9:\exists z\in\mathbb C^9,\
 f_j^{\mathbb Z}(e,z)=0\text{ for all }j\}.
\]
One exact complex pair \((e,z)\) would give a counterexample, since
the three star pivots already guarantee independence.

\begin{proposition}[Two certificates for the star ideal]
\label{prop:branch-certificates}
The membership \(1\in I_\star\) is equivalent to the existence of
\(a,h_j,k_j\in P\) and an integer \(N\ge1\) satisfying
\begin{align}
 1-aD &= \sum_{j=1}^{31}h_j f_j ,
       \label{eq:target-closed-certificate}\\
 D^N &= \sum_{j=1}^{31}k_j f_j .
       \label{eq:target-open-certificate}
\end{align}
\end{proposition}
\begin{proof}
The first identity makes \(D\) a unit in \(P/I_\star\), so it rules
out the collision branch \(D=0\). The second rules out the branch
where \(D\) is invertible. In \(P\) we have the explicit combination
\[
 1=\Bigl(\sum_{i=0}^{N-1}(aD)^i\Bigr)(1-aD)+a^N D^N .
\]
Substitution yields a certificate for one, including when the quotient
has nilpotents. Conversely, if \(1\in I_\star\), take \(a=0\) and
\(N=1\). Both right-hand representations then exist.
\end{proof}
Both polynomial identities remain to be found. The closed branch is
the first smaller calculation to try. Its kernel line gives a
17-variable chart covering it up to leaf permutations. A unit-ideal
calculation there would prove that branch empty. A certificate in the
original 18 variables could be replayed directly in
\eqref{eq:target-closed-certificate}.

The orbit count gives a second sufficient route. A nonempty geometric
fiber at five would contain at least eight free \(S_3\)-orbits, so
\(\dim_{\mathbb F_5}(P/I_\star)^{S_3}\le7\) would suffice. The missing
data on the whole chart are a span through \(1\) of at most seven
invariants and product identities for every span element against all
127 generators. The degree-seven dual calculation forces degree at
least eight in some product identity with \(1\). The existence of
such a certificate at degree eight or higher remains open.

\subsection{Selecting a witness from its encoded instance}
\label{sec:selector-target}

\paragraph{Uniform selector target for a specified verifier.}
Fix a polynomial-time binary verifier \(V(x,w)\) with a known,
polynomial-time computable padded witness length \(m(|x|)\) bounded by
a polynomial. Put \(W_x=\{w\in\{0,1\}^{m(|x|)}:V(x,w)=1\}\).
Construct a single deterministic polynomial-time algorithm which, on
every \(x\), returns an element of \(W_x\) if it is inhabited and
returns \(\bot\) exactly if \(W_x\) is empty. The explicit affine
family of Section~\ref{sec:disclosure} realizes this target by row
reduction. Requiring it for every such verifier is equivalent to
\(\PT=\NP\). A total selector decides feasibility. Conversely, under
\(\PT=\NP\) the prefix feasibility language below is decidable in
polynomial time for each fixed \(V\).

\paragraph{Two ways to discharge the target.}
A decomposition can supply a stratum index \(c(x)\), a representative
of its nonempty model fiber (or an empty marker), and a map taking each
valid model representative back to a valid original witness over \(x\).
It must also prove that every inhabited original fiber has an inhabited
model fiber at \(c(x)\). These conditions imply selector correctness
in Lean. To obtain the time bound, classification, representative
selection, and transport must all be uniform on the encoded inputs,
with polynomial bit cost including intermediate data. Hardt's theorem
provides product charts for a fixed semialgebraic projection
\cite[Theorem~4.1]{costeSemialgebraic}. For the projection onto
\(E_{\mathrm{lift}}\) above, the charts exist in real and imaginary
coordinates. A growing family of binary inputs needs its own effective
construction and bounds.

An alternative is to construct, in polynomial time from \((x,u)\), an
explicit polynomial-size binary affine system \(Q_{x,u}v=b_{x,u}\)
such that it is consistent exactly when some \(w\in W_x\) starts with
\(u\). First test the empty prefix. From a feasible prefix \(u\), take
\(u0\) when that extension is feasible and \(u1\) otherwise. After at
most \(m(|x|)+1\) consistency tests the procedure returns a valid
original witness or has certified emptiness. For the decoded affine
HCP family the construction just appends coordinate equations. For a
succinct verifier beyond this control, the proof obligation is to give
the prefix construction and prove the equivalence for every prefix.
